\documentclass[aspectratio=169]{beamer}

\usetheme{Madrid}
\usecolortheme{default}
\setbeamertemplate{navigation symbols}{}
\usepackage{eso-pic}
\usepackage[brazil]{babel}
\shorthandoff{"}
\usepackage[T1]{fontenc}
\usepackage[utf8]{inputenc}
\usepackage{lmodern}
 
\usepackage{tikz}
\usetikzlibrary{arrows.meta, positioning, shapes.geometric, fit, calc, backgrounds}
\usepackage{listings}
\usepackage{xcolor}

\definecolor{meuazul}{RGB}{0, 79, 151}
\setbeamercolor{structure}{fg=meuazul}

\title[Metodologias Ágeis]{Metodologias Ágeis e Planning Poker}
\subtitle{Engenharia da Qualidade e Confiabilidade}
\author{Sergio — UniRV}
\titlegraphic{\includegraphics[width=2cm]{logo.png}}

\begin{document}

% Slide de Título
\frame[plain]{\titlepage}

\AddToShipoutPictureFG{
    \put(\LenToUnit{.92\paperwidth},\LenToUnit{.185\paperheight})
    {\vtop{{\null}\makebox{\includegraphics[width=.8cm,keepaspectratio]{logo.png}}}}}

% --- PARTE 1: METODOLOGIAS ÁGEIS (10 Slides) ---

\begin{frame}{1) Contexto: Por que Ágil?}
  O modelo tradicional (\textbf{Cascata}) apresenta desafios em projetos de software:
  \begin{itemize}
    \item Feedback tardio (somente ao final).
    \item Dificuldade em lidar com mudanças de requisitos.
    \item Risco de entregar algo que o cliente não precisa mais.
  \end{itemize}
  \vspace{0.4em}
  A abordagem \textbf{Ágil} surge para focar em entregas rápidas e adaptabilidade.
\end{frame}

\begin{frame}{2) O Manifesto Ágil}
  Lançado em 2001, foca em 4 valores fundamentais:
  \begin{itemize}
    \item \textbf{Indivíduos e interações} mais que processos e ferramentas.
    \item \textbf{Software em funcionamento} mais que documentação abrangente.
    \item \textbf{Colaboração com o cliente} mais que negociação de contratos.
    \item \textbf{Responder a mudanças} mais que seguir um plano.
  \end{itemize}
\end{frame}

\begin{frame}{3) Os 12 Princípios Ágeis}
  Alguns pontos cruciais:
  \begin{itemize}
    \item Satisfação do cliente através de entrega contínua.
    \item Aceitar mudanças de requisitos, mesmo no fim do desenvolvimento.
    \item Entregar software funcionando com frequência (semanas ou meses).
    \item Cooperação diária entre pessoas de negócio e desenvolvedores.
    \item Atenção contínua à excelência técnica e bom design.
  \end{itemize}
\end{frame}

\begin{frame}{4) Ciclo de Vida Iterativo e Incremental}
  \centering
  \begin{tikzpicture}[node distance=2cm, auto]
    \node[draw, circle, meuazul, thick] (plan) {Planejar};
    \node[draw, circle, meuazul, thick, right=of plan] (dev) {Desenvolver};
    \node[draw, circle, meuazul, thick, below=of dev] (test) {Testar};
    \node[draw, circle, meuazul, thick, below=of plan] (review) {Revisar};
    
    \draw[->, thick, meuazul] (plan) -- (dev);
    \draw[->, thick, meuazul] (dev) -- (test);
    \draw[->, thick, meuazul] (test) -- (review);
    \draw[->, thick, meuazul] (review) -- (plan);
    
    \node[below=of review, yshift=1cm, text width=6cm, align=center] {Pequenas iterações que geram incrementos de software reais.};
  \end{tikzpicture}
\end{frame}

\begin{frame}{5) Principais Frameworks Ágeis}
  Existem diversas formas de aplicar a agilidade:
  \begin{itemize}
    \item \textbf{Scrum}: Focado em papéis, eventos e artefatos (o mais comum).
    \item \textbf{Kanban}: Focado em fluxo visual e limitação de trabalho em progresso (WIP).
    \item \textbf{XP (Extreme Programming)}: Focado em práticas de engenharia (TDD, Pair Programming).
  \end{itemize}
\end{frame}

\begin{frame}{6) Visão Geral do Scrum}
  O Scrum organiza o trabalho em:
  \begin{itemize}
    \item \textbf{Papéis}: Product Owner, Scrum Master, Developers.
    \item \textbf{Artefatos}: Product Backlog, Sprint Backlog, Increment (Pronto!).
    \item \textbf{Eventos}: Sprint, Sprint Planning, Daily Scrum, Review e Retrospective.
  \end{itemize}
\end{frame}

\begin{frame}{7) A Sprint: O Pulsação do Scrum}

\begin{itemize}
    \item Períodos de tempo fixo (\textit{time-box}), geralmente de 1 a 4 semanas.
    \item No início, planeja-se o que será feito (\textit{Planning}).
    \item No final, demonstra-se o valor entregue (\textit{Review}).
    \item O objetivo é sempre ter um incremento funcional ao final de cada Sprint.
\end{itemize}

\vspace{0.5cm}

\centering
\begin{tikzpicture}[node distance=2.5cm, every node/.style={font=\small}]

% Estilos
\tikzstyle{phase} = [rectangle, rounded corners, minimum width=2.5cm, minimum height=1cm, text centered, draw=blue!60, fill=blue!20]
\tikzstyle{arrow} = [thick,->,>=stealth]

% Nós
\node (planning) [phase] {Planning};
\node (dev) [phase, right of=planning, xshift=1.5cm] {Desenvolvimento};
\node (review) [phase, right of=dev, xshift=1.5cm] {Review};
\node (retro) [phase, below of=dev, yshift=-0.5cm] {Retrospectiva};

% Setas principais
\draw [arrow] (planning) -- (dev);
\draw [arrow] (dev) -- (review);

% Retorno (ciclo)
\draw [arrow] (review) |- (retro);
\draw [arrow] (retro) -| (planning);

\end{tikzpicture}

\vspace{0.3cm}
\small{Ciclo iterativo de uma Sprint no Scrum}

\end{frame}

\begin{frame}{8) Benefícios das Metodologias Ágeis}
  \begin{itemize}
    \item \textbf{Qualidade}: Testes frequentes e revisões constantes.
    \item \textbf{Visibilidade}: Transparência sobre o progresso real do projeto.
    \item \textbf{Mitigação de Risco}: Falhas são detectadas cedo.
    \item \textbf{Time to Market}: Entrega de valor para o cliente de forma mais rápida.
  \end{itemize}
\end{frame}

\begin{frame}{9) Mindset Ágil}
  Não é apenas sobre ferramentas, é sobre \textbf{cultura}:
  \begin{itemize}
    \item Foco no valor para o usuário.
    \item Transparência e Honestidade.
    \item Inspeção e Adaptação constante.
    \item Times auto-organizados e motivados.
  \end{itemize}
\end{frame}

\begin{frame}{10) Resumo da Agilidade}
  \begin{block}{Conclusão Parte 1}
    Ser ágil significa ser capaz de se mover rapidamente e com facilidade. No software, significa abraçar a incerteza e evoluir através da colaboração constante e entregas frequentes.
  \end{block}
\end{frame}

% --- PARTE 2: PLANNING POKER (5 Slides) ---

\begin{frame}{11) O que é Planning Poker?}

É uma técnica baseada em \textbf{consenso} para estimar o \textbf{esforço} ou tamanho relativo de tarefas (User Stories).

\begin{itemize}
    \item Evita o "efeito âncora".
    \item Envolve todo o time técnico na estimativa.
    \item Usa Story Points em vez de horas exatas.
\end{itemize}

\vspace{0.4cm}

\centering
\begin{tikzpicture}[every node/.style={font=\small}]

% Estilos
\tikzstyle{dev} = [circle, draw=blue!60, fill=blue!20, minimum size=1cm]
\tikzstyle{card} = [rectangle, rounded corners, draw=red!60, fill=red!20, minimum width=1cm, minimum height=1.2cm]
\tikzstyle{arrow} = [->, thick]

% Nós (desenvolvedores)
\node (d1) at (0,0) {Dev 1};
\node (d2) [right=2cm of d1] {Dev 2};
\node (d3) [right=2cm of d2] {Dev 3};

% Cartas
\node (c1) [card, below=0.8cm of d1] {3};
\node (c2) [card, below=0.8cm of d2] {5};
\node (c3) [card, below=0.8cm of d3] {8};

% Consenso
\node (consenso) [draw=green!60, fill=green!20, rounded corners, below=2.5cm of d2] {Discussão $\rightarrow$ Consenso};

% Setas
\draw[arrow] (c1) -- (consenso);
\draw[arrow] (c2) -- (consenso);
\draw[arrow] (c3) -- (consenso);

\end{tikzpicture}

\vspace{0.3cm}
\small{Cada membro escolhe uma carta; diferenças levam à discussão até consenso}

\end{frame}

\begin{frame}{12) O Baralho de Cartas}

As cartas geralmente seguem a \textbf{Sequência de Fibonacci} (modificada):

\begin{center}
\textbf{0, 1, 2, 3, 5, 8, 13, 21, 40, 100}
\end{center}

\begin{itemize}
    \item \textbf{Por que Fibonacci?} Porque reflete a incerteza: quanto maior a tarefa, mais difícil é ser preciso.
    \item \textbf{? (Interrogação)}: "Não entendi a tarefa".
    \item \textbf{Xícara de Café}: "Preciso de um intervalo".
\end{itemize}

\vspace{0.4cm}

\centering
\begin{tikzpicture}[every node/.style={font=\small}]

% Estilo das cartas
\tikzstyle{card} = [rectangle, rounded corners=0.2cm, draw=blue!60, fill=blue!20, minimum width=0.9cm, minimum height=1.2cm]

% Cartas Fibonacci
\node[card] (c0) at (0,0) {0};
\node[card] (c1) at (1.2,0) {1};
\node[card] (c2) at (2.4,0) {2};
\node[card] (c3) at (3.6,0) {3};
\node[card] (c5) at (4.8,0) {5};
\node[card] (c8) at (6.0,0) {8};
\node[card] (c13) at (7.2,0) {13};
\node[card] (c21) at (8.4,0) {21};
\node[card] (c40) at (9.6,0) {40};
\node[card] (c100) at (10.8,0) {100};

% Cartas especiais
\tikzstyle{special} = [rectangle, rounded corners=0.2cm, draw=red!60, fill=red!20, minimum width=1.2cm, minimum height=1.2cm]

\node[special] (q) at (3,-2) {?};
\node[special] (coffee) at (7,-2) {$\cup$};

% Legendas
\node at (3,-3) {\scriptsize Não entendi};
\node at (7,-3) {\scriptsize Pausa};

\end{tikzpicture}

\vspace{0.3cm}
\small{Cartas crescem em escala para representar maior incerteza e complexidade}

\end{frame}

 

\begin{frame}{14) Buscando o Consenso}
  E se houver divergência (ex: um deu 2 e outro deu 13)?
  \begin{itemize}
    \item Os extremos explicam seus pontos de vista.
    \item O time discute riscos ou complexidades que outros podem não ter visto.
    \item Uma nova rodada de votação é feita até que o time chegue a um acordo.
  \end{itemize}
\end{frame}

\begin{frame}{15) Vantagens do Planning Poker}
  \begin{itemize}
    \item \textbf{Engajamento}: Todo o time participa e compreende a tarefa.
    \item \textbf{Descoberta}: Discussões revelam requisitos ocultos ou dificuldades técnicas.
    \item \textbf{Precisão Relativa}: Estimar por comparação é mais natural do que estimar horas.
    \item \textbf{Divertido}: Torna o planejamento menos cansativo.
  \end{itemize}
\end{frame}

\begin{frame}
  \centering
  \Huge Obrigado!\\
  \vspace{1cm}
  \large Dúvidas?
\end{frame}

\end{document}
